\documentclass[12pt,a4paper]{article}
\usepackage[T2A]{fontenc}
\usepackage[utf8]{inputenc}
\usepackage[ukrainian]{babel}
\usepackage[left=2.5cm,right=2cm,top=2cm,bottom=2.5cm]{geometry}
\usepackage{graphicx}
\usepackage{float}
\usepackage{booktabs}
\usepackage{array}
\usepackage{amsmath}
\usepackage[labelsep=colon]{caption}
\usepackage{enumitem}
\usepackage{fancyvrb}
\usepackage{indentfirst}
\usepackage{tikz}
\usetikzlibrary{positioning,arrows.meta,shapes.geometric}

\graphicspath{{figures/}}
\fvset{fontsize=\small,frame=single,framesep=2mm}
\setlist[enumerate]{leftmargin=*}
\setlength{\parskip}{0.3em}
\setlength{\emergencystretch}{3em}

% Титульна сторінка: #1 — номер лабораторної роботи
\newcommand{\labtitlepage}[1]{%
\begin{titlepage}
\begin{center}
\large
Міністерство освіти і науки України\\
Національний університет ``Львівська політехніка''\\
Кафедра прикладної математики
\vspace*{\fill}

{\LARGE Звіт}\\[0.5em]
{\Large про виконання лабораторної роботи №#1}\\[0.5em]
{\Large з курсу «Основи інженерії даних»}
\vspace*{\fill}
\end{center}
\begin{flushright}
\begin{minipage}{0.4\textwidth}
\large
Виконала:\\
студентка групи ПМІД-11\\
Мазур Д. М.\\[1.5em]
Перевірив:\\
Алєксєєв В. І.
\end{minipage}
\end{flushright}
\vspace*{\fill}
\begin{center}
\large Львів — 2026
\end{center}
\end{titlepage}
\setcounter{page}{2}}

\newcommand{\rsec}[1]{\par\medskip\noindent\textbf{#1}\ }

\begin{document}
\labtitlepage{0}

\rsec{Тема:} Підготовча робота. Налаштування середовища розробки та робота з реляційною базою даних засобами мови Python.

\rsec{Мета роботи:} Відновити у пам'яті знання щодо роботи з базами даних та написання програмного коду. Нагадати навички розгортання середовища для розробки, включно з підключенням системи контролю версій.

\rsec{Завдання до роботи:}
\begin{enumerate}
  \item Встановити СУБД PostgreSQL.
  \item Встановити та налаштувати середовище для програмування мовою Python.
  \item Підключити систему контролю версій (GitHub).
  \item Створити просту БД та наповнити її згенерованими даними.
  \item Написати програмний код мовою Python для виконання запитів на вибірку, зміну та видалення даних у створеній БД.
\end{enumerate}

\rsec{Хід виконання роботи:}
\begin{enumerate}
  \item \textbf{Встановлення СУБД PostgreSQL.} PostgreSQL версії 15 було встановлено через пакетний менеджер Homebrew та запущено як фонову службу. Для лабораторної роботи створено окрему базу даних \texttt{lab0\_db}:
\begin{Verbatim}
brew install postgresql@15
brew services start postgresql@15
createdb lab0_db
\end{Verbatim}

  \item \textbf{Налаштування середовища Python.} Використано інтерпретатор Python 3.9. Для ізоляції залежностей проєкту створено віртуальне середовище \texttt{venv}, у яке встановлено драйвер PostgreSQL \texttt{psycopg2-binary} (2.9) та бібліотеку генерації тестових даних \texttt{Faker} (37.12):
\begin{Verbatim}
python3 -m venv venv
source venv/bin/activate
pip install psycopg2-binary faker
\end{Verbatim}

  \item \textbf{Підключення системи контролю версій.} Для всіх лабораторних робіт курсу створено один Git-репозиторій \texttt{fundam\_of\_data\_engineering}, який опубліковано на GitHub (\texttt{github.com/daria-mazur/fundam\_of\_data\_engineering}). Кожна лабораторна робота розміщена в окремій папці (\texttt{lab0\_db}, \dots, \texttt{lab4\_db}) і фіксується окремим комітом. Віртуальне середовище, кеш Python і великі файли даних виключено з репозиторію через \texttt{.gitignore}.

  \item \textbf{Створення бази даних.} Скрипт \texttt{lab0\_db/main.py} створює таблицю \texttt{students} (табл.~\ref{tab:students}). Команда \texttt{CREATE TABLE IF NOT EXISTS} робить створення ідемпотентним, тому скрипт можна запускати повторно.

\begin{table}[H]
\centering
\caption{Структура таблиці \texttt{students}}
\label{tab:students}
\begin{tabular}{lll}
\toprule
Поле & Тип & Обмеження / опис \\
\midrule
\texttt{id} & \texttt{SERIAL} & первинний ключ, автоінкремент \\
\texttt{full\_name} & \texttt{VARCHAR(100)} & ПІБ студента, \texttt{NOT NULL} \\
\texttt{email} & \texttt{VARCHAR(100)} & \texttt{UNIQUE NOT NULL} \\
\texttt{gpa} & \texttt{NUMERIC(5,2)} & середній бал (60--100) \\
\texttt{created\_at} & \texttt{TIMESTAMP} & за замовчуванням \texttt{CURRENT\_TIMESTAMP} \\
\bottomrule
\end{tabular}
\end{table}

  \item \textbf{Наповнення згенерованими даними.} Функція \texttt{seed\_data()} генерує записи за допомогою \texttt{Faker} з українською локаллю (\texttt{uk\_UA}): ПІБ, email та випадковий середній бал. Усі запити параметризовані (\texttt{\%s}), що захищає від SQL-ін'єкцій. Конструкція \texttt{ON CONFLICT (email) DO NOTHING} пропускає рядки з email, який уже є в таблиці, тому повторний запуск скрипта не завершується помилкою порушення унікальності.
\begin{Verbatim}
insert_query = """
INSERT INTO students (full_name, email, gpa)
VALUES (%s, %s, %s)
ON CONFLICT (email) DO NOTHING;
"""
for _ in range(count):
    cursor.execute(insert_query,
                   (fake.name(), fake.unique.email(),
                    round(random.uniform(60.0, 100.0), 2)))
    inserted += cursor.rowcount
conn.commit()
\end{Verbatim}

  \item \textbf{Запити на вибірку, зміну та видалення.} Для кожної операції написано окрему функцію: \texttt{fetch\_all\_students()} (\texttt{SELECT \dots ORDER BY id}), \texttt{update\_student\_gpa()} (\texttt{UPDATE \dots WHERE id = \%s}) та \texttt{delete\_student()} (\texttt{DELETE \dots WHERE id = \%s}). Кожна функція відкриває з'єднання, виконує запит і фіксує транзакцію через \texttt{commit()}. Для \texttt{UPDATE} і \texttt{DELETE} виводиться \texttt{cursor.rowcount}, тобто фактична кількість змінених рядків.

  \item \textbf{Демонстрація роботи.} Сценарій у блоці \texttt{\_\_main\_\_} створює таблицю, додає 5 записів, виводить вміст таблиці, оновлює середній бал першого студента до 99.50, видаляє цього студента та щоразу показує результат. Фрагмент виводу програми:
\begin{Verbatim}
[+] Таблицю 'students' успішно створено/перевірено.
[+] Додано 5 згенерованих записів.

--- Список студентів у БД ---
ID: 4 | ПІБ: Ярослав Самойленко | Email: andriishcherban@... | GPA: 75.62
ID: 5 | ПІБ: Максим Тичина | Email: nadiiamalyshko@... | GPA: 65.92
...
[+] Оновлено GPA для студента з ID=4 на 99.5 (змінено рядків: 1).
ID: 4 | ПІБ: Ярослав Самойленко | Email: andriishcherban@... | GPA: 99.50
...
[-] Видалено студента з ID=4 (видалено рядків: 1).
\end{Verbatim}
  Після видалення запис з ID=4 відсутній у вибірці, а значення \texttt{rowcount = 1} підтверджує, що кожна операція змінила рівно один рядок.
\end{enumerate}

\rsec{Висновки:} У ході виконання лабораторної роботи №0 було налаштовано середовище для подальших робіт курсу: СУБД PostgreSQL 15, Python 3.9 з віртуальним середовищем \texttt{venv} та Git-репозиторій на GitHub. Створено базу даних \texttt{lab0\_db} з таблицею \texttt{students} і наповнено її згенерованими україномовними даними за допомогою бібліотеки Faker. Реалізовано всі операції CRUD через драйвер \texttt{psycopg2} з параметризованими запитами та явним керуванням транзакціями. Обробка конфлікту унікального ключа (\texttt{ON CONFLICT DO NOTHING}) та ідемпотентне створення таблиці дозволяють запускати скрипт повторно без помилок.

\end{document}
